\MFQLead{对冲失效与研究材料答案}
\begin{enumerate}
  \item 还剩 Gamma、Vega、Theta、高阶项、参数误差、跳跃、离散调仓、成本和成交误差。Delta 只是当前模型下的一阶局部导数。
  \item 单路径无法估计 bias、方差或尾部；研究者还可能事后挑路径。至少应固定路径生成协议并报告全体误差分布。
  \item 令现金 $B_i$ 先增长为 $B_{i-1}e^{r\Delta t}$，再令 $q_i=\Delta_i-\Delta_{i-1}$，则 $B_i=B_{i-1}e^{r\Delta t}-q_iS_i-c|q_i|S_i$。
  \item bias 揭示平均方向，RMSE 同时惩罚偏差与波动，分位数不假设对称并揭示尾部。三者不能互相替代。
  \item 事先确定 24 步结果无成本 RMSE 为 $1.416360$，成本后为 $1.434142$；12 与 52 步应在同一随机路径和成本规则下重算，不能只比较其中最优者。
  \item 在某一步把 $S$ 乘以跳跃因子并保持旧 Delta 到下一个调仓点。报告跳跃方向、时点、大小、成本以及误差分位数变化。
  \item 错误报告只给路径编号、没有全体路径和预先选择规则。修复应事先确定种子与路径数，提交所有统计或可重建账本，而不是再挑一条看似典型的路径。
  \item 合格 Capstone 包含合约与报价时间、曲线/股息、模型假设、四种定价误差、曲面约束、Greek 验证、多路径对冲、摩擦压力、故障注入、限制与一键命令。
\end{enumerate}

\section*{逐题推导与核对}
\begin{enumerate}[leftmargin=*]
  \item Delta 中性只使一阶价格变动项 $\Delta\,\mathrm dS$ 在局部消失。二阶展开还留下 $\frac12\Gamma(\mathrm dS)^2$，参数变化留下 Vega，时间衰减留下 Theta；离散调仓、跳跃、报价误差、成交失败和融资成本又会产生模型外误差。因此应按风险来源分解，而不是把 Delta=0 写成“无风险”。
  \item 设路径误差为 $e_j=\Pi_T^{(j)}-H^{(j)}$。一条路径只能给出一个 $e_j$，不能估计 $\mathbb E e_j$、$\mathbb E e_j^2$ 或尾分位数；若路径是事后挑选的，甚至会有选择偏差。至少冻结路径生成器、种子、调仓规则和成本口径，再报告均值、标准差、P5/P50/P95、最大损失与有效样本数。
  \item 先让上一期现金按无风险利率增长：$B_i^-=B_{i-1}e^{r\Delta t}$。若持仓从 $\Delta_{i-1}$ 改为 $\Delta_i$，交易量为 $q_i=\Delta_i-\Delta_{i-1}$，比例成本为 $cS_i|q_i|$，则
  \[
  B_i=B_i^- -q_iS_i-cS_i|q_i|,
  \qquad V_i=B_i+\Delta_iS_i.
  \]
  这个顺序把成交发生在 $t_i$ 写清楚；若把成本从总组合价值中重复扣除，或让未成交数量凭空消失，就破坏自融资账本。
  \item $\operatorname{bias}=\mathbb E[e]$ 只描述方向；$\operatorname{RMSE}=\sqrt{\mathbb E[e^2]}$ 同时包含系统偏差和随机波动；分位数 $Q_{0.05},Q_{0.50},Q_{0.95}$ 描述非对称尾部。关系 $\mathrm{RMSE}^2=\operatorname{Var}(e)+\operatorname{bias}^2$ 说明即使 bias 为零，尾部仍可能很大。
  \item 固定同一批路径，分别取 $N=12,24,52$ 个调仓点。对每个 $N$ 计算无成本与成本后 $e_j$，再用同一组路径得到均值、RMSE 和分位数；已知 $N=24$ 的 RMSE 为 $1.416360$、成本后为 $1.434142$。如果只换路径，误差变化同时混入 Monte Carlo 噪声，不能归因于调仓频率。
  \item 在预先指定的时刻令 $S_{t_k+}=J S_{t_k-}$，并保持上一次 Delta 到下一次调仓。逐行记录跳跃前后价格、理论持仓、实际成交量、成本和 PnL；比较 $J>1$ 与 $J<1$ 的尾部。跳跃通常放大 Gamma 误差，成本可能因价格水平改变而同时变化，故必须分列报告。
  \item “完美复制”路径只说明一个样本点的 $e_j$ 接近零；它没有给出路径抽样协议、置信区间、尾部、成本和模型错设敏感性。修复时先冻结全体路径，再自动生成统计表；禁止按结果挑路径，保留失败路径编号和原因。
  \item Capstone 应按顺序提交：报价快照与单位；模型和参数假设；解析/树/PDE/Monte Carlo 定价互核；Delta/Gamma/Vega 检查；多路径对冲分布；成本、跳跃和波动率错设压力；失败 fixture；最后给限制和一键执行入口。每一步都保存输入哈希和独立 oracle，才能从报价追到最终意见。
\end{enumerate}
